\chapter{Project Risk Assessment \& Management}

\section{Project Risk Assessment}

• Risk assessment consists of brainstorming to identify risks that might
affect a project. Team members informally select the top two or three
risks that present the greatest threat to the project and develop a plan
to manage them.

interpersonal conflict - 3, 1 
overleaf failure - 4, 4
bottleneck - 4, 5
miscommunication - 5, 5
poor time estimation - 3, 2
external/internal advice delay - 2, 3 
team member absence - 1, 1




\textcolor{blue}{Project Insert risk map here}
\begin{table}[h]
    \centering
        \caption{Project Risk Matrix}
    \begin{tabular}{|c|c|c|c|c|c|c|}
         &  \\
         & 
    \end{tabular}
    \label{tab:placeholder}
\end{table}

\section{Risk Management Plans}
• Risk management plans should incorporate both actions that can be
taken to reduce the likelihood of failure (i.e., preventive measures),
and actions that can be taken in the event of failure (i.e., contingency
plans)